\chapter{Примеры решений}

\section{П1. Перевод делением и обратная проверка}\label{ex:p1}
\textbf{Условие.} Перевести $181_{10}$ в основания 2, 8 и 16.

Последовательные деления на 2:
\begin{center}
\begin{tabular}{r|r|r}
Делимое & Частное & Остаток\\\hline
181 & 90 & 1\\90 & 45 & 0\\45 & 22 & 1\\22 & 11 & 0\\
11 & 5 & 1\\5 & 2 & 1\\2 & 1 & 0\\1 & 0 & 1
\end{tabular}
\end{center}
Читаем остатки снизу вверх: $10110101_2$.

Для основания 8: $181=22\cdot8+5$, $22=2\cdot8+6$, $2=0\cdot8+2$. Ответ: $265_8$.

Для основания 16: $181=11\cdot16+5$, $11=0\cdot16+11$. Остаток 11 обозначаем B, поэтому ответ $B5_{16}$.

\textbf{Проверка:} $128+32+16+4+1=181$; $2\cdot64+6\cdot8+5=181$; $11\cdot16+5=181$.

\textbf{Частая ошибка:} читать остатки сверху вниз. Получившийся ответ обязательно проверьте весами разрядов.

\section{П2. Перевод группировкой битов}\label{ex:p2}
\textbf{Условие.} Перевести $1101011011_2$ в основания 4, 8 и 16.
\begin{itemize}
    \item По 2 бита: \texttt{11 01 01 10 11}, значит $31123_4$.
    \item По 3 бита: \texttt{001 101 011 011}, значит $1533_8$.
    \item По 4 бита: \texttt{0011 0101 1011}, значит $35B_{16}$.
\end{itemize}
\textbf{Проверка:} $35B_{16}=3\cdot256+5\cdot16+11=859$;
$1533_8=512+5\cdot64+3\cdot8+3=859$;
$31123_4=3\cdot256+64+16+2\cdot4+3=859$.

\textbf{Частая ошибка:} дополнить нулями справа. Это умножает число на степень двойки. Дополнение целой записи делается слева.

\section{П3. Четыре арифметических действия}\label{ex:p3}
\textbf{Сложение:} $1101_2+1011_2$. Справа налево:
$1+1=10_2$; $0+1+1=10_2$; $1+0+1=10_2$; $1+1+1=11_2$.
Получаем $11000_2$. Проверка: $13+11=24$.

\textbf{Вычитание:} $B2_{16}-6D_{16}$. Для младшего разряда занимаем единицу у B: старшая цифра становится A, младшее значение --- $16+2=18$. Тогда $18-13=5$, $10-6=4$. Ответ $45_{16}$; проверка $178-109=69$.

\textbf{Умножение:} $26_8\cdot3_8$. Младший разряд: $6\cdot3=18_{10}=22_8$, пишем 2, переносим 2. Старший: $2\cdot3+2=8_{10}=10_8$, пишем 0, переносим 1. Ответ $102_8$; проверка $22\cdot3=66$.

\textbf{Деление:} $247_8:5_8$. Первые две цифры дают $24_8$. Подбираем $4_8$, поскольку $5_8\cdot4_8=24_8$. Вычитаем, сносим 7: следующая цифра частного 1, остаток 2. Ответ: частное $41_8$, остаток $2_8$.
Проверка в десятичной системе: $167=5\cdot33+2$, $2<5$.

\textbf{Частая ошибка:} переносить десяток, как в десятичной записи. В каждом разряде переносится полная группа из $p$ единиц.

\section{П4. Диапазон и минимальная ширина}\label{ex:p4}
\textbf{Условие.} Найти диапазоны 10-битного unsigned и дополнительного кода. Выбрать минимальную знаковую ширину для $[-70,90]$.

$2^{10}=1024$, поэтому unsigned: $[0,1023]$; дополнительный код: $[-512,511]$. В обоих случаях 1024 значения.

Для 7 битов знаковый диапазон $[-64,63]$ --- не подходит ни одна из требуемых границ. Для 8 битов $[-128,127]$ --- подходят обе. Ответ: 8 битов.

\textbf{Частая ошибка:} считать, что раз есть отдельный знак, диапазон дополнительного кода симметричен. Отрицательных значений в нём на одно больше, чем положительных.

\section{П5. Число в четырёх кодах}\label{ex:p5}
\textbf{Условие.} Представить $-29$ в 8 битах; для смещённого кода взять bias $=128$.

Положительное $29=16+8+4+1$ имеет запись \texttt{00011101}.
\begin{itemize}
    \item Прямой: ставим знак 1 перед семибитным модулем. Получаем \texttt{10011101} ($9D_{16}$).
    \item Обратный: инверсия \texttt{00011101} даёт \texttt{11100010} ($E2_{16}$).
    \item Дополнительный: добавляем 1 к обратному коду, получаем \texttt{11100011} ($E3_{16}$). Проверка: $256-29=227$.
    \item Со смещением: $-29+128=99=64+32+2+1$. Код \texttt{01100011} ($63_{16}$).
\end{itemize}
\textbf{Частая ошибка:} инвертировать только биты модуля. При переходе к обратному и дополнительному кодам инвертируется вся $n$-битная запись положительного числа.

\section{П6. Интерпретация и расширение}\label{ex:p6}
\textbf{Условие.} Прочитать слово \texttt{11100011} как unsigned, дополнительный код и код со смещением 128. Расширить первые две интерпретации до 16 битов.

Беззнаковое значение: $128+64+32+2+1=227$.
В дополнительном коде старший бит равен 1, поэтому $227-256=-29$.
В коде со смещением получаем $227-128=99$.

Расширение unsigned: \texttt{00000000 11100011}, то есть $00E3_{16}$ и число 227.
Знаковое расширение: \texttt{11111111 11100011}, то есть $FFE3_{16}$ и число $65507-65536=-29$.

\textbf{Частая ошибка:} всегда добавлять нули. Для отрицательного числа в дополнительном коде это изменяет значение.

\section{П7. Сложение и независимость флагов}\label{ex:p7}
\textbf{Условие.} Выполнить в 8 битах $70+65$, $(-70)+(-65)$ и сложение кодов $C8_{16}+38_{16}$.

\begin{enumerate}
    \item $70=46_{16}$, $65=41_{16}$. Сумма $87_{16}$, то есть \texttt{10000111}. Беззнаково это 135, знаково $-121$. CF=0, поскольку $135<256$; OF=1, поскольку точное $135>127$. ZF=0, SF=1.
    \item $-70$ кодируется $BA_{16}$, $-65$ --- $BF_{16}$. Сумма кодов $179_{16}$, сохраняем $79_{16}$, то есть 121. CF=1. Точная знаковая сумма $-135<-128$, поэтому OF=1. ZF=0, SF=0.
    \item $C8_{16}+38_{16}=100_{16}$. Сохраняем \texttt{00000000}; CF=1, ZF=1, SF=0. Знаковые операнды $-56$ и 56 дают ровно 0, поэтому OF=0.
\end{enumerate}
\textbf{Частая ошибка:} определять OF по наличию девятого бита. Девятый бит суммы отвечает за CF; OF проверяется по знаковой арифметике.

\section{П8. Вычитание и флаг займа}\label{ex:p8}
\textbf{Условие.} Выполнить 8-битное $42-75$ через дополнительный код и указать флаги SUB.

$42=\texttt{00101010}$; $75=\texttt{01001011}$. Инвертируем второй код и добавляем 1:
\[
\texttt{10110100}+1=\texttt{10110101}.
\]
Складываем:
\[
\texttt{00101010}+\texttt{10110101}=\texttt{11011111}.
\]
Ответ $DF_{16}$, знаковое значение $223-256=-33$, совпадающее с $42-75$.

\textbf{Флаги SUB:} CF=1, так как $42<75$ и требуется заём; OF=0, так как $-33$ представимо; ZF=0; SF=1.

У вспомогательного сложения кодов $42+181=223$ переноса нет. CF сложения и CF исходного вычитания определяются для разных операций.

\section{П9. Побитовые операции и маски}\label{ex:p9}
\textbf{Условие.} Даны $x=96_{16}$, $y=3A_{16}$, ширина 8 битов.
\begin{center}
\begin{tabular}{l|c|c}
Операция & Биты результата & Hex\\\hline
$x$ & \texttt{10010110} & 96\\
$y$ & \texttt{00111010} & 3A\\
AND & \texttt{00010010} & 12\\
OR & \texttt{10111110} & BE\\
XOR & \texttt{10101100} & AC\\
NOT $x$ & \texttt{01101001} & 69
\end{tabular}
\end{center}
Например, единицы AND остаются только в позициях 4 и 1, где они есть в обоих операндах.

Независимые изменения исходного $x$:
\begin{itemize}
    \item Установить бит 5: маска $20_{16}$, $96_{16}\mathbin{\mathrm{OR}}20_{16}=B6_{16}$.
    \item Сбросить бит 7: маска сохранения $7F_{16}$, результат $16_{16}$.
    \item Инвертировать бит 1: XOR с $02_{16}$, результат $94_{16}$.
    \item Проверить бит 4: AND с $10_{16}$ даёт $10_{16}\ne0$, значит бит установлен.
\end{itemize}
\textbf{Частая ошибка:} вместо сброса применять XOR. XOR меняет 1 на 0, но одновременно меняет 0 на 1; сброс должен оставлять исходный нулевой бит нулевым.

\section{П10. Сдвиги одного и того же слова}\label{ex:p10}
\textbf{Условие.} Дано 8-битное слово $B5_{16}=\texttt{10110101}$.
Выполнить независимо сдвиги на 2 бита.

\begin{center}
\begin{tabular}{l|c|c}
Сдвиг & Результат & Hex\\\hline
Логический влево & \texttt{11010100} & D4\\
Логический вправо & \texttt{00101101} & 2D\\
Арифметический вправо & \texttt{11101101} & ED\\
Циклический влево & \texttt{11010110} & D6
\end{tabular}
\end{center}
При левом логическом сдвиге старшие \texttt{10} теряются, справа появляются \texttt{00}. При циклическом те же \texttt{10} возвращаются справа.

Исходное знаковое число равно $181-256=-75$. Арифметический сдвиг даёт $\lfloor-75/4\rfloor=-19$, код $ED_{16}$. Деление $-75/4$ для целых в C даёт $-18$, поскольку округляется к нулю.

\textbf{Проверка левого сдвига:} $181\cdot4=724$, $724\bmod256=212=D4_{16}$. Он не сохранил точное произведение.
